% Aporte de Angello Mendoza Bermello. Incluir desde el informe principal.
\subsection{Implementación con gRPC y Protocol Buffers}
\label{sec:grpc}

Con gRPC el trabajo cambia de lugar. En TCP armábamos los bytes y los leíamos a
mano; aquí declaramos una interfaz y llamamos a una función que vive en otro
proceso. Del resto se encarga la biblioteca: empaqueta en binario, transporta por
HTTP/2 y vuelve a armar el mensaje al otro lado~\cite{thomson2022http2}.

\subsubsection{El contrato}

Todo empieza en \path{proto/lamport.proto}, un archivo de texto donde se declaran
los mensajes y las funciones que se pueden llamar. Lo esencial es esto:

\begin{verbatim}
message Mensaje {
  string sender = 1;
  int64  timestamp = 2;
  string message = 3;
  string run_id = 4;
  string message_id = 5;
  string kind = 6;
}

message Acuse {
  string sender = 1;
  int64  timestamp = 2;
  string message = 3;
  string run_id = 4;
  string message_id = 5;
  string kind = 6;
  bool   ok = 7;
}

service NodoLamport {
  rpc EnviarMensaje (Mensaje) returns (Acuse);
}

service ControlNodo {
  rpc Preparar (Vacio) returns (EstadoListo);
  rpc Intercambiar (Orden) returns (Resultado);
}
\end{verbatim}

\texttt{NodoLamport} es el servicio del enunciado y es la única llamada que
cronometramos.

\texttt{ControlNodo} lo agregamos nosotros, para que el lanzador pueda pedirle a un
nodo que envíe un mensaje. Así el envío lo origina el mismo proceso que lleva el
reloj, y cada nodo conserva un solo contador.

Los campos son los del JSON de TCP, más tres que nos hacían falta:
\texttt{run\_id} identifica la corrida, \texttt{message\_id} empareja cada
solicitud con su acuse y \texttt{kind} distingue una de otra. El sello viaja como
entero de 64 bits, no como texto.

\subsubsection{Del contrato al programa}

El código de Python se genera con \texttt{grpcio-tools}, mediante el script
\path{src/grpc/generar_stubs.py}. Los archivos generados quedan en el repositorio:
no hay que regenerarlos para ejecutar el proyecto.

Aquí cada nodo es un proceso, no un hilo como en TCP. Cada uno escucha en su
puerto (5000, 5001 y 5002) y atiende los dos servicios en el mismo puerto. Cuando
le llega una solicitud, aplica la regla de recepción del reloj, anota el evento y
responde con su propio sello. Al detenerlo con Ctrl+C escribe su CSV.

Para la prueba concurrente, el lanzador sincroniza los tres envíos con una
barrera y cada nodo los atiende con un grupo de hilos. Las tres llamadas quedan en
vuelo al mismo tiempo.

\subsubsection{El reloj y lo que se mide}

Cada nodo tiene un reloj y dos papeles: atiende lo que le llega y origina lo que
le ordenan. Usamos el reloj lógico de Lamport~\cite{lamport1978}, que ordena
dos eventos cuando hay un mensaje de por medio, pero no cuando no lo hay. Las dos
reglas están en el mismo bloque y bajo el mismo candado:

{\small
\begin{verbatim}
def sellar_envio(self, peer, run_id, message_id, kind, carga):
    with self._candado:
        self._reloj += 1
        self._anotar("send", peer, run_id, message_id, kind, carga)
        return self._reloj

def registrar_recibo(self, peer, run_id, message_id, kind, carga, sello_remoto):
    with self._candado:
        self._reloj = max(self._reloj, int(sello_remoto)) + 1
        self._anotar("receive", peer, run_id, message_id, kind, carga)
        return self._reloj
\end{verbatim}
}

El RTT se toma dentro del proceso que envía, con \texttt{time.perf\_counter()},
desde antes de sellar el mensaje hasta que vuelve el acuse. Fuera de la medición
quedan el arranque de los procesos y el establecimiento de los canales. Por eso el
protocolo descarta una corrida de calentamiento: en las primeras operaciones el
rendimiento todavía no es estable~\cite{traini2023steady}. Dentro quedan la
serialización, el transporte y la decodificación, igual que en el cliente TCP.

\subsubsection{Errores y costo}

El contrato tipado ahorra trabajo. Los campos tienen tipo y el código generado
construye el objeto, así que no hay que analizar el JSON a mano. El servidor
revisa el emisor y los campos obligatorios, y responde con un estado
(\texttt{INVALID\_ARGUMENT}) en vez de quedarse callado. El cliente comprueba que
el acuse corresponda a lo que envió. Si un nodo está apagado, la llamada termina
con \texttt{UNAVAILABLE}, y eso separa un problema de red de uno de datos.

Nada de esto sale gratis. Debajo hay HTTP/2, la biblioteca de gRPC y una
codificación binaria, y eso se nota en el tiempo. Los estudios que comparan gRPC
con otros protocolos en microservicios encuentran resultados que dependen del
escenario: en uno gRPC queda primero en tiempo de respuesta~\cite{niswar2024microservices}
y en otro las diferencias no llegan a ser significativas~\cite{ain2025grpc}.
Nuestras cifras van en esa misma línea y se discuten más adelante.
